% BEGIN SUPPLEMENTARY MATERIAL: supplementary_now (1).tex -- Multi-View Silhouette Refinement -- implementation details
\chapter{Silhouette Refinement Implementation Details}
\label{app:refinement}
This appendix specifies the test-time procedure used with the silhouette
objective in Sec.~\ref{supp:silhouette_refinement}.

\paragraph{Parameter selection.}
% TODO(author): specify the validation split or recordings used for refinement
% parameter selection.
Refinement hyperparameters were selected on a held-out validation split and
kept fixed for test evaluation.

\paragraph{Rendering and search.}
We render the full evaluation mesh at half image resolution with pixel stride
two, restricting rendering to a crop around the projected object padded by
the bounding-box diagonal.
Starting from the network pose, coordinate descent uses rotation steps
\(4^\circ,2^\circ,1^\circ,0.5^\circ,\) and \(0.25^\circ\).
We evaluate \(16\) rotation initializations and accept the refined pose only
when its silhouette objective improves over the network initialization by a
margin of \(3.0\).
% TODO(author, supplementary integration): confirm how the 3.0 acceptance
% margin is scaled relative to the summed fractional objective in
% Sec.~\ref{supp:silhouette_refinement}; the source does not specify a scaling.

\paragraph{Object-specific degrees of freedom.}
For the drill, refinement optimizes rotation only.
For the screwdriver, we additionally scan rotation about the object-frame
\(z\)-axis over \(\pm25^\circ\), using a \(2.5^\circ\) coarse step followed
by a \(0.5^\circ\) fine step.
Translation refinement with a \(4\,\mathrm{mm}\) step is enabled for the
screwdriver with predicted masks, but not with benchmark masks.

\paragraph{Computation.}
On \(12\) CPU workers, refinement of the complete test set requires
approximately one hour per object.
% END SUPPLEMENTARY MATERIAL: Multi-View Silhouette Refinement -- implementation details
